\chapter{El problema financiero}\label{cap:problema_financiero}

Este capítulo plantea en lenguaje físico, y sin dar por sabido el
cálculo estocástico, el problema que resuelve el trabajo: qué contratos
se valoran, por qué su precio está determinado y por qué las opciones
que dependen de la trayectoria piden otro método. Las demostraciones se
citan \cite{tfgm, shreve, hull}.

\section{Contratos y pagos}\label{sec:contratos}

\subsection{Opciones europeas}\label{subsec:europeas}

Consideramos un mercado con dos activos. El primero es un activo con
riesgo ---una acción, por ejemplo---, cuyo precio $S_t$ fluctúa de forma
aleatoria; en los capítulos siguientes hará el papel de la coordenada
de una partícula. El segundo es una cuenta bancaria sin riesgo, cuyo
saldo crece de forma determinista, $B_t = B_0\, e^{rt}$, con tipo de
interés $r$ constante. Un euro cobrado en $T$ equivale, por tanto, a
$e^{-r\tau}$ euros en el instante $t$, con $\tau = T - t$: el factor de
descuento convierte cantidades de una fecha a otra, como un cambio de
unidades, y solo tiene sentido comparar precios referidos a la misma
fecha.

Un \emph{derivado} es un contrato cuyo pago depende del precio de otro
activo, llamado \emph{subyacente}. El más simple es la \emph{opción de
compra europea} (\emph{call}) con precio de ejercicio (\emph{strike})
$K$ y vencimiento $T$: da a su titular el derecho, pero no la
obligación, de comprar una unidad del subyacente en el instante $T$ al
precio $K$. Si $S_T > K$, el titular ejerce y gana $S_T - K$; si
$S_T \leq K$, no ejerce y no gana nada. El \emph{pago} (\emph{payoff})
en $T$ es
\begin{equation}\label{eq:pago_call}
  g(S_T) \;=\; (S_T - K)^+ \;\eqdef\; \max(S_T - K,\, 0).
\end{equation}
La opción de venta europea (\emph{put}) paga $(K - S_T)^+$, y en general
llamamos europea a cualquier opción que paga una función $g(S_T)$ del
precio en el vencimiento. Su rasgo estructural es que el pago depende
solo del estado final: en lenguaje físico, $g$ es un observable que se
mide en $t = T$ sobre la posición $S_T$, y la historia que llevó hasta
allí es irrelevante. Para valorarla basta conocer la ley de una sola
variable aleatoria, $S_T$.

El problema es determinar cuánto vale hoy, en $t < T$, el derecho a
cobrar $g(S_T)$ en $T$. Conviene adelantar por qué hace falta un modelo
del subyacente. Si el pago fuera lineal, $S_T - K$ (un contrato a plazo
o \emph{forward}), su precio quedaría fijado sin hipótesis alguna sobre
la dinámica de $S$: comprar hoy una unidad del subyacente y pedir
prestados $K e^{-r\tau}$ reproduce ese pago en cualquier escenario, así
que el contrato vale $S_t - K e^{-r\tau}$. La \emph{call}, en cambio,
tiene un codo en $S_T = K$: es convexa, y su valor depende de toda la
distribución de $S_T$ y no solo de su media. Por eso entra la
volatilidad del subyacente y por eso hace falta un modelo
(sec.~\ref{sec:modelo_bs}).

Las opciones que estudia este trabajo pagan, en cambio, una
\emph{funcional de la trayectoria}: el pago en $T$ es
$\Phi[S(\cdot)]$, con $S(\cdot) = \{S_t\}_{0 \leq t \leq T}$, y la
europea es el caso particular $\Phi[S(\cdot)] = g(S_T)$. Las dos
familias del trabajo son las asiáticas, que dependen de una integral a
lo largo de la trayectoria, y las barrera, que dependen de sus valores
extremos.

\subsection{Opciones asiáticas}
\label{subsec:asian_concepto}\label{sec:asian_definicion}
\label{subsec:asian_promedios}

Una opción asiática es un derivado cuyo pago depende del valor
\emph{promedio} del subyacente sobre un periodo de monitorización, en
lugar de su valor terminal: el promedio es más difícil de manipular que
una observación aislada y, como tiene menos varianza, la opción es más
barata que la europea equivalente. Sobre el periodo $[0, T]$, los dos promedios canónicos
son el aritmético y el geométrico:
\begin{align}
  A_T &\;\eqdef\; \frac{1}{T}\int_0^T S_t\, dt,
  \label{eq:promedio_aritmetico}\\
  G_T &\;\eqdef\; \exp\!\left(\frac{1}{T}\int_0^T \log S_t\, dt\right).
  \label{eq:promedio_geometrico}
\end{align}
El caso canónico, y el que se valora en este trabajo, es la \emph{call}
de precio de ejercicio fijo (\emph{fixed-strike call}), que paga
$(A_T - K)^+$ o $(G_T - K)^+$ según el promedio; la \emph{put} y las
variantes de precio de ejercicio flotante, que comparan $S_T$ con el
promedio, se tratan del mismo modo. En la práctica el promedio se toma
sobre un número finito de fechas; aquí se usa siempre la
monitorización continua, que es su límite.

\medskip

\noindent\textbf{Linealidad en el log-precio.}
La observación que estructura los Caps.~\ref{cap:fuente_pared}
y~\ref{cap:espectro}\enprep{} aparece al pasar al log-precio $X_t = \log S_t$.
El logaritmo del promedio geométrico es lineal en la trayectoria; el
del aritmético, no:
\begin{align}
  \log G_T &\;=\; \frac{1}{T}\int_0^T X_t\, dt,
  \label{eq:log_geom_lineal}\\
  \log A_T &\;=\; \log\!\left(\frac{1}{T}\int_0^T e^{X_t}\, dt\right).
  \label{eq:log_arit_no_lineal}
\end{align}
En el modelo de Black--Scholes, $X_t$ es un proceso gaussiano
(sec.~\ref{sec:modelo_bs}). Una combinación lineal de variables
gaussianas es gaussiana, así que $G_T$ es lognormal y su opción admite
un tratamiento elemental; $A_T$, integral de exponenciales de un
proceso gaussiano, no es lognormal ni tiene una densidad cerrada
elemental, y es esta no linealidad la que hace el caso aritmético
fundamentalmente más difícil. Por la desigualdad entre las medias
aritmética y geométrica, consecuencia de la convexidad de la
exponencial, $G_T \leq A_T$ en cada trayectoria; para volatilidades y
plazos realistas la diferencia es apreciable, y el caso aritmético no
se reduce al geométrico.

\subsection{Opciones barrera}
\label{subsec:barrera_unica_payoff}\label{sec:barrera_definicion}
\label{subsec:doble_barrera_otros_tipos}

Una opción barrera paga como una europea, pero solo si el subyacente ha
cruzado, o no, un nivel $B$ fijado en el contrato. La opción se
desactiva (\emph{knock-out}) o se activa (\emph{knock-in}) al tocar la
barrera, que puede estar por debajo (\emph{down}) o por encima
(\emph{up}) del precio inicial. El caso canónico es la
\emph{down-and-out call}, con $B < S_0$:
\begin{equation}\label{eq:payoff_DO_call}
  C_{DO}(T) \;=\; (S_T - K)^+ \cdot
  \1\!\left\{\min_{0 \leq t \leq T} S_t > B\right\}.
\end{equation}
El indicador anula el pago si el subyacente tocó $B$ en algún instante,
sea cual sea $S_T$: la trayectoria entra a través de su mínimo. Como
toda trayectoria toca la barrera o no la toca,
$\1\{\min_t S_t > B\} + \1\{\min_t S_t \leq B\} = 1$, y los pagos de
la \emph{down-and-out} y la \emph{down-and-in} con la misma barrera
suman el de la europea. Lo mismo ocurre con sus precios, porque un
precio que no fuera aditivo permitiría una ganancia sin riesgo
(\emph{paridad in--out}):
\begin{equation}\label{eq:paridad_in_out}
  C_{DO}(t, S) + C_{DI}(t, S) \;=\; C_{\text{vanilla}}(t, S).
\end{equation}
Basta, pues, valorar las opciones \emph{out}; la europea se obtiene en
el Cap.~\ref{cap:bs_operadores}\enprep. En la opción de \emph{doble barrera}
se fijan dos niveles $B_d < B_u$, y el pago se cobra solo si el
subyacente permanece en $(B_d, B_u)$ durante toda la vida del contrato:
la \emph{double knock-out call} paga
$(S_T - K)^+\, \1\{B_d < \min_t S_t,\ \max_t S_t < B_u\}$. En todo el trabajo la barrera tiene nivel constante, se monitoriza de
forma continua y no paga compensación (\emph{rebate}) al desactivarse.

\section{No arbitraje y replicación en un modelo de un periodo}\label{sec:no_arbitraje}

El principio que fija los precios es el de \emph{no arbitraje}: no
existe ninguna cartera que cueste cero o menos hoy, no pierda en ningún
escenario y gane en alguno con probabilidad positiva. Es una condición
de equilibrio: si tal cartera existiera, todos la explotarían y los
precios se ajustarían hasta eliminarla. Veamos su efecto en el modelo
más simple posible \cite[sec.~2.2]{tfgm}.

\medskip

\noindent\textbf{El modelo.}
Hay un único periodo, de duración $\Delta t$. El subyacente pasa de
$S_0$ a $S_0 u$ con probabilidad $p$ o a $S_0 d$ con probabilidad
$1 - p$, con $0 < d < u$, y la cuenta bancaria multiplica su saldo por
$e^{r\Delta t}$. El no arbitraje exige
\begin{equation}\label{eq:no_arbitraje_un_periodo}
  d \;<\; e^{r\Delta t} \;<\; u .
\end{equation}
Si $e^{r\Delta t} \geq u$, vender el activo en corto (venderlo prestado
y devolverlo al final) y depositar lo obtenido nunca pierde y gana si el
activo baja; si $e^{r\Delta t} \leq d$, pedir prestado para comprar el
activo nunca pierde y gana si sube.

\medskip

\noindent\textbf{La réplica.}
Una opción que vence al final del periodo paga $V_u$ si el activo sube
y $V_d$ si baja; para la \emph{call}, $V_u = (S_0 u - K)^+$ y
$V_d = (S_0 d - K)^+$. Formamos una cartera con $\Delta$ unidades del
subyacente y $\vartheta$ euros en la cuenta, y exigimos que pague lo mismo
que la opción en los dos estados:
$\Delta\, S_0 u + \vartheta\, e^{r\Delta t} = V_u$ y
$\Delta\, S_0 d + \vartheta\, e^{r\Delta t} = V_d$. Este sistema lineal
tiene solución única,
\begin{equation}\label{eq:delta_un_periodo}
  \Delta \;=\; \frac{V_u - V_d}{S_0 u - S_0 d},
  \qquad
  \vartheta \;=\; e^{-r\Delta t}\,\frac{u V_d - d V_u}{u - d}.
\end{equation}
La cartera cuesta hoy $\Delta S_0 + \vartheta$. Si la opción se negociara a
otro precio, comprar la más barata de las dos y vender la más cara daría
una ganancia segura sin inversión, así que el no arbitraje fija
\begin{equation}\label{eq:precio_un_periodo}
  V_0 \;=\; \Delta S_0 + \vartheta
  \;=\; e^{-r\Delta t}\bigl[\,q\, V_u + (1 - q)\, V_d\,\bigr],
  \qquad
  q \;\eqdef\; \frac{e^{r\Delta t} - d}{u - d}.
\end{equation}

\noindent\textbf{Lectura.}
La probabilidad $p$ no aparece en \eqref{eq:precio_un_periodo}. La razón
es geométrica: con dos estados posibles, los pagos al final del periodo
forman un espacio vectorial de dimensión dos, y el subyacente y la
cuenta, con pagos $(S_0 u, S_0 d)$ y $(e^{r\Delta t}, e^{r\Delta t})$
linealmente independientes, son una base. Cualquier opción es un vector
de ese espacio, y su precio queda fijado por linealidad al descomponerla
en la base, sin que intervenga lo probable que sea cada estado. En
particular, el precio no depende de cuánto se espera que crezca el
subyacente: en el modelo continuo, no dependerá de su rendimiento
esperado $\mu$. Además, $\Delta$ en \eqref{eq:delta_un_periodo} es un
cociente incremental, $\Delta V / \Delta S$, versión discreta de la
cobertura $\Delta = \partial_S V$.
Con tres estados la réplica deja de ser posible en general (mercado
\emph{incompleto}, \cite[sec.~2.2.5]{tfgm}); el modelo de
Black--Scholes, límite $\Delta t \to 0$ de un modelo de muchos periodos
\cite[secs.~2.3--2.4]{tfgm}, es completo.

\section{La probabilidad neutral al riesgo}\label{sec:neutral_riesgo}

La condición \eqref{eq:no_arbitraje_un_periodo} equivale a $0 < q < 1$,
así que $q$ puede leerse como una probabilidad. Llamamos $\Q$ a la que
asigna probabilidad $q$ a la subida y $1 - q$ a la bajada. Con ella,
\eqref{eq:precio_un_periodo} dice que el precio de la opción es su pago
descontado promediado con $\Q$,
\begin{equation}\label{eq:precio_q_un_periodo}
  V_0 \;=\; e^{-r\Delta t}\, \E_\Q[V_1],
\end{equation}
y un cálculo directo muestra qué caracteriza a $\Q$:
\begin{equation}\label{eq:q_crecimiento_r}
  \E_\Q[S_1] \;=\; q\, S_0 u + (1 - q)\, S_0 d \;=\; S_0\, e^{r\Delta t}.
\end{equation}
Bajo $\Q$ el activo con riesgo crece en media al tipo sin riesgo, como
la cuenta bancaria: de ahí el nombre de probabilidad \emph{neutral al
riesgo}. Dicho de otro modo, el precio descontado $e^{-rt}S_t$ tiene al
final del periodo la misma esperanza que al principio. Un proceso cuyo
valor esperado futuro, condicionado a la información presente, coincide
con su valor presente es una \emph{martingala}: un juego justo, sin
deriva que explotar. Por \eqref{eq:precio_q_un_periodo}, también el
precio descontado de la opción es una martingala bajo $\Q$.

Dos precisiones evitan malentendidos. La primera: $\Q$ no es una
predicción de lo que hará el mercado. Las probabilidades reales están
en $p$, que no interviene; $\Q$ son los únicos pesos con los que
promediar los pagos es compatible con el no arbitraje, y usarlos es un
recurso de cálculo. La segunda: $\Q$ y la probabilidad real $\PP$
asignan probabilidad cero a los mismos sucesos (aquí a ninguno, porque
$0 < p, q < 1$). Cambian los pesos de los escenarios, no cuáles son
posibles; en esto consiste que sean \emph{equivalentes}.

El resultado se generaliza en el \emph{teorema fundamental de la
valoración} \cite[sec.~5.4]{shreve}: un mercado no admite arbitraje si
y solo si existe una probabilidad $\Q$ equivalente a $\PP$ bajo la cual
los precios descontados de todos los activos son martingalas, y esa
$\Q$ es única si y solo si el mercado es completo. El precio de un
derivado replicable es entonces su pago descontado promediado con $\Q$;
en el modelo binomial, la construcción explícita está en
\cite[secs.~2.2--2.3]{tfgm}.

\section{El modelo de Black--Scholes}\label{sec:modelo_bs}

El modelo de Black--Scholes describe el subyacente con las hipótesis más
simples compatibles con un precio positivo y aleatorio: volatilidad
$\sigma$ y tipo $r$ constantes, trayectorias sin saltos y negociación
continua y sin costes. Su ingrediente es el movimiento browniano, del
que aquí basta lo mínimo; su física (el paseo aleatorio, la ecuación del
calor y la ecuación de Langevin) se desarrolla en el
Cap.~\ref{cap:fundamentos}\enprep.

\medskip

\noindent\textbf{El browniano y la regla $dW^2 = dt$.}
El movimiento browniano $W_t$ empieza en $W_0 = 0$, tiene trayectorias
continuas e incrementos independientes en intervalos disjuntos, con
\begin{equation}\label{eq:incremento_browniano}
  W_{t + \Delta t} - W_t \;\sim\; N(0, \Delta t).
\end{equation}
Un incremento típico mide $\sqrt{\Delta t}$, mucho más que $\Delta t$
cuando $\Delta t$ es pequeño: las trayectorias son continuas pero no
derivables en ningún punto, y los desarrollos de Taylor deben llevarse a
segundo orden en $\Delta W$. Ese término de segundo orden se comporta de
forma notable. Cada $(\Delta W)^2$ tiene media $\Delta t$ y varianza
$2\Delta t^2$; sumados sobre los $t/\Delta t$ pasos de $[0, t]$, dan una
variable de media $t$ y varianza $2t\,\Delta t$, que tiende a cero. En
el límite, la suma de los cuadrados de los incrementos deja de ser
aleatoria y vale $t$ en cada trayectoria (es la \emph{variación
cuadrática}, \cite[sec.~3.2.3]{tfgm}). En un desarrollo infinitesimal
esto se resume en la regla
\begin{equation}\label{eq:regla_dW2}
  (dW)^2 = dt, \qquad dW\, dt = 0, \qquad (dt)^2 = 0,
\end{equation}
donde las dos últimas significan que esos términos son de orden
superior a $dt$. Es la única regla nueva que necesita este capítulo.

\medskip

\noindent\textbf{El modelo.}
El precio del subyacente sigue la ecuación
\begin{equation}\label{eq:bs_eds_P}
  dS_t \;=\; \mu\, S_t\, dt + \sigma\, S_t\, dW_t,
\end{equation}
límite de $S_{t+\Delta t} = S_t\,(1 + \mu\,\Delta t + \sigma\,\Delta W)$:
el rendimiento relativo $dS/S$ tiene una deriva $\mu$ y un ruido de
amplitud $\sigma$, la \emph{volatilidad}. Para el log-precio
$X_t = \log S_t$, el desarrollo de Taylor a segundo orden, con
$(dS/S)^2 = \sigma^2\,dt$ por \eqref{eq:regla_dW2}, da
\begin{equation}\label{eq:logprecio_P}
  dX_t \;=\; \frac{dS_t}{S_t} - \frac{1}{2}\left(\frac{dS_t}{S_t}\right)^{\!2}
  \;=\; \Bigl(\mu - \frac{\sigma^2}{2}\Bigr)\, dt + \sigma\, dW_t .
\end{equation}
Es la ecuación de Langevin de una partícula browniana sometida a una
fuerza constante, con deriva $\mu - \sigma^2/2$ y coeficiente de difusión
$\sigma^2/2$. Sus coeficientes son constantes, así que
$X_t = X_0 + (\mu - \sigma^2/2)\,t + \sigma W_t$ es gaussiano y $S_t$ es
lognormal. El término $-\sigma^2/2$ procede de la concavidad del
logaritmo: subir un 10\% y bajar después un 10\% deja el precio
en $0{,}99$ veces el inicial. La volatilidad frena así el crecimiento
típico de una trayectoria, medido por la mediana
$S_0\,e^{(\mu - \sigma^2/2)t}$, aunque la media siga siendo
$S_0\,e^{\mu t}$. El mismo desarrollo, aplicado a una función
$f(t, S)$, da la \emph{regla de Itô} \cite[sec.~4.3]{tfgm}:
\begin{equation}\label{eq:regla_ito}
  df \;=\; \Bigl(\partial_t f + \mu S\, \partial_S f
  + \frac{1}{2}\sigma^2 S^2\, \partial_S^2 f\Bigr)\, dt
  + \sigma S\, \partial_S f\, dW_t ,
\end{equation}
cuyo único rasgo no clásico es el término en $\partial_S^2 f$, procedente
de $(dS)^2 = \sigma^2 S^2\,dt$.

\medskip

\noindent\textbf{Bajo $\Q$.}
El teorema de Girsanov \cite[sec.~4.4]{tfgm} construye la probabilidad
neutral al riesgo del modelo continuo: bajo $\Q$, el proceso
$\widetilde W_t = W_t + \frac{\mu - r}{\sigma}\,t$ es un browniano, y
\begin{equation}\label{eq:bs_eds_Q}
  dS_t \;=\; r\, S_t\, dt + \sigma\, S_t\, d\widetilde W_t,
  \qquad
  dX_t \;=\; \Bigl(r - \frac{\sigma^2}{2}\Bigr)\, dt
  + \sigma\, d\widetilde W_t,
\end{equation}
cuya segunda ecuación es el punto de partida del
Cap.~\ref{cap:bs_operadores}\enprep. Como en \eqref{eq:q_crecimiento_r}, el
subyacente crece en media al tipo $r$. La deriva cambia, pero la
volatilidad no puede hacerlo, por la razón que da la regla
\eqref{eq:regla_dW2}: la suma de los cuadrados de los incrementos de $X$
en $[0, t]$ vale $\sigma^2 t$ con probabilidad uno, y dos probabilidades
equivalentes coinciden en los sucesos de probabilidad uno. La
volatilidad se lee en una sola trayectoria; la deriva no puede
determinarse a partir de una trayectoria de duración finita, y por eso
es lo único que el cambio de probabilidad modifica. En el
Cap.~\ref{cap:bs_operadores}\enprep, multiplicar por una exponencial del
log-precio elimina la deriva del hamiltoniano: es una reponderación de
trayectorias del mismo tipo, vista como transformación gauge.

\seccionesprevistas{2.5--2.6}
